\section{Introduction}

Construction work entails substantial safety and health risks. The U.S. Bureau of Labor Statistics reported 1,032 fatal occupational injuries in construction and extraction occupations in 2024; common events included falls, slips, trips, and transportation incidents involving construction vehicles or machinery \cite{bls2024}. Beyond traumatic injuries, work with silica-containing materials such as concrete and mortar can expose workers to high concentrations of respirable crystalline silica over long periods, increasing the risks of silicosis, lung cancer, and chronic obstructive pulmonary disease \cite{niosh2024}. Concrete floor leveling and troweling further combine a large work area, repeated passes, and a narrow time window governed by concrete setting. Manual work therefore requires sustained physical exertion and repetitive motion. Field physiological measurements of concrete slab placing and finishing crews showed that most manual workers experienced high energy expenditure and cardiopulmonary demand, with a greater risk of physical fatigue than equipment operators \cite{abdelhamid1999}. Work quality also depends on operator experience, proficiency, and moment-to-moment condition \cite{albers2004,kobayashi1996}. Releasing workers from hazardous, repetitive, and physically demanding floor operations while improving the consistency of the finished surface is therefore a problem of direct engineering and social value for construction robotics.

Floor-troweling robots provide a practical route toward mechanization. Early field trials showed that a remotely operated troweling robot could perform multiple operations from wood floating to steel troweling. It treated approximately 1,000~m$^2$ across five sites and achieved an average productivity of about 317~m$^2$/h under the reported test conditions \cite{kobayashi1996}. For an order-of-magnitude comparison, the average production rate of manual steel troweling is approximately 13.9~m$^2$/h \cite{methvin2026}; the reported robotic rate is therefore about 22.8 times the rate of one manual worker. Subsequent work introduced electric actuation, laser localization, task planning, and remote monitoring to ride-on power trowels, allowing them to follow specified paths within a predefined floor area \cite{hwangbo1999}.

Most existing troweling robots use a disc mounted beneath or in front of the base. Because the tool remains fixed relative to the base, the base path largely determines the tool path, and back-and-forth routes can be planned with a fixed working width \cite{kobayashi1996,hwangbo1999}. Swing-arm configurations have recently attracted attention because they enlarge the end-disc workspace without requiring a large base. Such a robot can place its circular disc beside the base or into locally constrained regions while the base remains clear of walls, columns, and openings, thereby partially decoupling the traversable base path from the tool trajectory. The arm can also increase the effective coverage width. For a fixed-disc platform, the nominal width of a straight pass is mainly limited by the disc diameter. A swing-arm platform can instead move the disc laterally while the base advances, so that the effective width depends on the arm length, angular range, and disc diameter and can exceed the disc diameter \cite{lin2025}. At the same base speed, provided that the swing frequency keeps consecutive sweeps contiguous without introducing gaps, the theoretical covered area per unit time can increase from approximately $v(2R)$ for a fixed disc to $v w_{\mathrm{eff}}$ for a swing arm, where $w_{\mathrm{eff}}>2R$. Previous coverage-planning research likewise shows that repeated coverage, average velocity, and the smoothness of turns jointly affect task time and energy consumption \cite{galceran2013,lee2011}. The swing-arm configuration may also reduce unnecessary driving and sharp turns on wet concrete and provide more admissible poses near corners, narrow-passage entrances, and obstacles \cite{lin2025}.

\begin{figure}[htbp]
  \centering
  \includegraphics[width=0.82\linewidth]{figure_1_swing_arm_troweling_robot_photo.jpg}
  \caption{A swing-arm troweling robot operating on site. The tracked base, swing arm and actuator, and circular end disc are visible. Image source: Shangguan News; photograph provided by China Construction Eighth Engineering Division.}
  \label{fig:robot-photo}
\end{figure}

Mechanization and teleoperation, however, do not by themselves provide autonomy. A human teleoperator must continuously observe the floor and machine state while coordinating translation, steering, and troweling. In spaces divided by walls, columns, embedded components, and openings, occluded views and fragmented passages increase the burden of global route selection and can cause backtracking, detours, missed regions, and repeated work. Experiments with teleoperated mobile robots also show that manual driving time increases with environmental difficulty and depends on operator skill and the level of obstacle-avoidance assistance \cite{storms2016}. Existing swing-arm systems commonly rely on human--robot cooperation or predefined trajectories and do not yet generate complete coverage routes automatically for complex construction maps \cite{lin2025}. Continuous teleoperation also requires a skilled operator to remain on duty, adding training and on-site supervision costs. Empirical studies of construction-robot adoption further identify economic factors, work culture, and uncertain business returns as barriers to deployment \cite{ocs2024,delgado2019}. Without autonomous coverage planning, these operational burdens remain even after the troweling process is mechanized.

This task is an instance of coverage path planning (CPP): the tool must sweep the target area while avoiding obstacles and controlling missed coverage, repeated coverage, transfer distance, and execution time \cite{galceran2013}. Existing approaches include back-and-forth sweeps after cellular decomposition, grid or spanning-tree traversal, and strip generation with curvature-constrained connections for fixed-width implements. These approaches respectively address polygonal or occupancy-grid maps, regular regions, and applications such as agriculture in which the tool footprint is fixed \cite{choset1998,gabriely2001,bahnemann2021,polygoncoverage,mier2023,fields2cover,nobleo}. Mobile-manipulator methods can plan local motions for a base and an end tool, while SatAOI can identify the workable area of a swing-arm troweling robot in an obstacle map; neither, however, provides a complete continuous coverage route for a complex two-dimensional floor \cite{lin2025,thakar2022}. Section~\ref{sec:related-work} gives a detailed classification and discusses these boundaries.

CPP for a swing-arm troweling robot introduces three interdependent requirements. First, the planner must represent a complex free space containing openings, dense obstacles, and narrow connections. Second, at a fixed base pose, it must distinguish the disc positions and whole-robot collision states associated with different arm angles. Third, while the base moves continuously, actual coverage must be determined from arm angular velocity and the continuous disc sweep rather than from a static reachability envelope. Simply dilating the base path into a constant-width strip cannot satisfy all three requirements. Consequently, prior work lacks a complete planning chain from a construction map through decomposition and visitation order to coordinated base--arm motion and end-tool coverage evaluation.

To close this gap, this paper presents a hierarchical CPP method for known two-dimensional environments with obstacles. Bidirectional rectangular candidates and connectivity completion represent the complex free space. A feasible transition graph then organizes global visitation, and first visits and repeated transit visits are instantiated as local swing-assisted coverage motions and transfer motions, respectively. Finally, the coverage ratio and execution time are computed from the region actually swept by the circular disc along the continuous trajectory. The same robot model is thus used for map representation, arm behavior, and trajectory evaluation.

The remainder of the paper is organized as follows. Section~\ref{sec:related-work} reviews the development of concrete-troweling robots and the challenges posed by obstacle-rich construction environments, compares representative CPP methods, and summarizes the unresolved research gap. Section~\ref{sec:methodology} presents the robot and environment models, rectangular-candidate decomposition, connectivity completion, global visitation, and swing-assisted coverage within each rectangle. Section~\ref{sec:experiments} evaluates the method through a baseline comparison, candidate-scan strategy experiments, and a speed-sensitivity study. Section~\ref{sec:conclusion} summarizes the experimentally supported contributions, limitations, and future work.

\section{Related Work}
\label{sec:related-work}

\subsection{Concrete Troweling Robots and Application Challenges}

After concrete has been placed and initially leveled, it must be troweled repeatedly before setting and hardening. Troweling removes local height variations and rough texture, redistributes the surface paste, and produces a flat, dense finish through compaction and polishing. The operation covers a large area and is highly repetitive, while its start time, speed, and number of passes are constrained by the state of concrete setting. Starting too early can disturb an unstable surface; starting too late can prevent the desired flatness and finish. Site practice has therefore long depended on the experience and continuous involvement of skilled workers \cite{albers2004,kobayashi1996}. A troweling robot moves the disc and treats the surface at a relatively stable speed and cadence, reducing physical workload and the effect of worker-state variability on quality consistency \cite{kobayashi1996,shin2003}. Its role is thus not merely to replace a hand tool, but to convert a process that requires continuous manual operation into one that a mechanical system can execute consistently.

Existing troweling robots can be grouped into three broad configurations according to the position of the disc relative to the base. In the first, the circular disc is mounted under or near the center of the base. This compact and mechanically stable arrangement is suitable for fixed-width operation on open floors, but the tool footprint essentially translates with the base path. In the second, the disc is mounted in front of the base, making the forward work area easy for an operator or sensor to observe, but requiring the base to approach the treatment boundary and offering limited lateral reach. In the third, the disc is mounted at the end of a swing arm. Controlled relative motion between the base and disc then allows the robot to cover lateral regions, areas near obstacles, and narrow-passage entrances while the base remains at a safe distance. This paper focuses on the third configuration for two reasons. The arm can actively reposition the disc relative to the base, making coverage near walls, columns, openings, and narrow passages more flexible. In addition, at the same base speed and with a contiguous sweep, the effective coverage width can exceed the disc diameter, increasing the area treated per unit time \cite{lin2025,galceran2013,lee2011}. Both advantages arise from coordinating base translation with arm swinging and cannot be captured by a conventional fixed-width strip model.

\begin{figure}[htbp]
  \centering
  \includegraphics[width=0.92\linewidth]{figure_2_disc_vehicle_relationships.png}
  \caption{Three arrangements of the troweling disc relative to the base: (a) disc beneath or near the center of the base; (b) front-mounted disc; and (c) disc at the end of a swing arm. Arrows indicate the direction of base travel.}
  \label{fig:disc-configurations}
\end{figure}

The swing-arm arrangement enlarges the platform workspace by mounting the disc at the end of a rotating long arm. Unlike a disc fixed beneath or in front of the base, the arm can treat lateral regions while the base maintains clearance and can create multiple base--tool configurations by changing its angle. This kinematic redundancy aids coverage near walls, columns, and narrow-passage entrances and may reduce base travel and sharp turns on wet concrete. The additional degree of freedom also makes the robot state more complex: $x_c$ and $y_c$ locate the center of the base in the ground frame, $\theta$ is the base heading, and $\phi$ is the arm angle relative to the base. Together, the four variables determine the poses of the base, arm, and disc in the construction plane. A region traversable by the base is not necessarily reachable by the disc, and a disc-reachable pose is not necessarily collision-free for the complete mechanism. Planning must therefore check the geometry of all three components and coordinate the base and arm throughout continuous motion \cite{lin2025}. The capability gained from the mechanical arrangement consequently depends on whether the control system can organize and execute the associated coordinated motion.

The evolution of concrete-troweling equipment can be understood along two axes: the disc--base relationship and the mode of control. Walk-behind and ride-on power trowels first mechanized the operation, but a worker still walked with or rode on the machine. Because the disc remained fixed beneath or in front of the body, the machine path effectively determined the disc path \cite{hwangbo1999}. Remote-control robots subsequently moved the operator outside the work area. Kobayashi et al. fixed front and rear troweling mechanisms to a lightweight mobile chassis and used wireless teleoperation and taught-motion replay to perform successive operations from wood floating to steel troweling. The system reduced worker exertion, but the troweling mechanisms still moved rigidly with the chassis, and coverage remained sensitive to operator experience, floor friction, and steering behavior \cite{kobayashi1996}.

Research then shifted from teleoperation toward autonomous travel in regular areas. Hwangbo et al. retained the fixed twin-rotor troweling mechanism of a ride-on machine, electrified controls such as steering, blade pitch, and acceleration, and added laser localization, reflective landmarks, and a construction-drawing-based task planner. The resulting machine could work autonomously along a prescribed pattern or planned route \cite{hwangbo1999}. This improved control autonomy without changing the rigid coupling between disc and body. More recent swing-arm robots alter that geometry: a tracked base provides locomotion, while a circular disc on an actively rotating arm moves independently relative to the base, partially decoupling traversable base routes from tool trajectories and enlarging the workable region near obstacles \cite{lin2025}. The evolution of troweling robots therefore encompasses both a progression from manual control through wireless teleoperation to partial autonomy and a transition from a body-fixed disc to an actively positioned one. The latter improves flexibility and efficiency in complex environments but introduces a new coordinated base--arm planning problem.

Among existing systems, the swing-arm platform considered by SatAOI is closest to the target of this paper. Mechanically, it retains active relative motion between the end disc and the base and can use multiple base--arm poses near walls, columns, openings, and narrow passages. At the planning level, SatAOI begins with an obstacle grid, explicitly models the base, arm, and disc and their collision rules, and performs a global search to identify mutually connected regions that the robot can enter and in which the disc can operate \cite{lin2025}. It therefore combines the spatial flexibility of a swing arm with environment-model-based autonomous decision making and provides a more complete foundation for automated troweling in complex construction scenes than a fixed tool or a purely teleoperated platform.

SatAOI nevertheless outputs a workable area rather than a complete motion route that covers it. It answers where the robot can work, but not in what order the regions should be visited, how the base should transfer between them, or how the base and arm should coordinate within each region. Key steps---area decomposition, global visitation, local coverage, and speed coordination---are still missing between AOI delimitation and an executable CPP solution. This is the direct starting point of the present study.

\subsection{Coverage Path Planning Methods}

Existing CPP methods can be compared along two dimensions: map-representation capability and fidelity of robot-behavior modeling. The first describes whether a planner can represent irregular boundaries, multiple obstacles, holes, and narrow connections and use them to establish adjacency and feasible transfers between coverage regions. The second describes whether the planner captures the physical dimensions, relative poses, motion modes, and collision constraints of the robot and tool rather than replacing the operation with a point, disc, or fixed-width strip. These dimensions define four broad capability classes. Idealized coverage methods address regular spaces with simplified robot models and mainly establish fundamental coverage mechanisms. Map-enhanced methods represent complex maps but retain a fixed tool footprint. Platform-constrained methods describe vehicle or manipulator motion constraints but primarily target regular or local workspaces. Environment--platform coupled methods preserve both complex environmental structure and realistic robot behavior.

Swing-arm troweling CPP requires both capabilities of the fourth class. Walls, columns, reserved openings, and narrow corridors divide the floor into irregular regions with complex connectivity, which the planner must represent and visit continuously. Meanwhile, disc position is jointly determined by base motion and active arm motion; coverage and collision status vary continuously with base heading and arm angle and cannot be approximated by a fixed footprint. A complete, continuous, and executable route therefore requires the construction map and the coordinated base--arm--disc behavior to be represented in the same planning model.

The first class comprises classical coverage algorithms that express the core coverage idea clearly but generally close the problem only for idealized maps and simplified robot models. Representative methods include Boustrophedon Cellular Decomposition (BCD) and Spanning Tree Coverage (STC). BCD partitions an area where scan-line connectivity changes and generates back-and-forth sweeps within each cell \cite{choset1998}; STC builds a spanning tree over a free-space grid and systematically traverses its boundary \cite{gabriely2001}. These methods establish the respective roles of cellular decomposition and graph traversal, but their coverage primitive is typically a point robot, disc robot, or fixed-width tool. Applied directly to swing-arm troweling, they cannot distinguish the relative poses of the base, arm, and disc or represent how arm-angle changes affect the instantaneous footprint and collision risk.

The second class offers stronger complex-map representation and includes reproducible open-source frameworks, but still models the robot and tool mainly through a fixed footprint. Examples include ETH Zurich's \texttt{polygon\_coverage\_planning}, Fields2Cover for agricultural vehicles, and the Nobleo Full Coverage Path Planner for ROS occupancy grids. The first handles general polygons with holes, uses BCD or trapezoidal decomposition to generate cells, constructs multiple sweep candidates for each cell, and performs global selection and collision-free transfer through an Equality Generalized Traveling Salesman Problem and a visibility graph \cite{bahnemann2021,polygoncoverage}. Fields2Cover supports non-convex fields, obstacles, headlands, swaths, route ordering, and curvature-constrained Dubins or Reeds--Shepp connections \cite{mier2023,fields2cover}. Nobleo's planner implements a Backtracking Spiral Algorithm/Spiral-STC-type method on a ROS occupancy grid and permits separate robot and tool radii \cite{nobleo}. These frameworks show that complex boundaries, obstacles, and engineering interfaces can be addressed systematically. Their tool width, implement width, or tool disc nevertheless remains essentially fixed relative to the base and cannot represent the time-varying sweep caused by continuous changes in $\phi$.

The third class emphasizes robot motion constraints, reachability, and collision relationships, but usually does not target complete coverage of a complex two-dimensional floor. It includes the nonholonomic vehicle connections in Fields2Cover and mobile-manipulator surface-coverage methods for tasks such as spray disinfection. Fields2Cover incorporates minimum turning radius and vehicle kinematics \cite{mier2023,fields2cover}, but targets agricultural swaths with a fixed implement width. Mobile-manipulator studies explicitly handle the base, arm, end tool, reachability, and occlusion. For example, spray-based surface disinfection can generate end-effector paths from three-dimensional surfaces and jointly select base positions and arm motions \cite{thakar2022}. Such methods describe how a robot executes a task in greater detail, but generally target regular sites, agricultural strips, or three-dimensional object surfaces rather than jointly performing obstacle handling, area decomposition, global traversal, and continuous ground-coverage evaluation on a known two-dimensional construction map.

Research has thus advanced in both complex-map representation and robot-behavior modeling, but the two have not yet been unified for swing-arm floor coverage. CPP frameworks with strong map support usually assume a fixed tool footprint and cannot describe the dynamic coverage created by a rotating arm. Methods that model vehicle kinematics, manipulator reachability, and component collisions in detail generally focus on regular sites, local operations, or three-dimensional surfaces and do not fully decompose and continuously traverse a complex two-dimensional construction map. SatAOI considers both an obstacle map and the geometry of the base, arm, and disc, but stops at delimiting the workable area and does not produce a visitation order, transfer route, or coordinated coverage trajectory within each region \cite{lin2025}. The central gap is therefore a swing-arm troweling CPP method that combines complex-map processing with high-fidelity robot-behavior modeling and generates a complete, continuous, executable route directly from the construction map.

Table~\ref{tab:cpp-comparison} compares representative approaches along these two dimensions. Work spanning more than one class is placed according to its principal contribution: Fields2Cover is treated as map-enhanced and SatAOI as platform-constrained. The stated limitations concern direct application to a mobile base with an active arm and circular disc; they are not general judgments of the methods or software.

\begin{longtable}{@{}>{\raggedright\arraybackslash}p{0.16\linewidth}>{\raggedright\arraybackslash}p{0.18\linewidth}>{\raggedright\arraybackslash}p{0.21\linewidth}>{\raggedright\arraybackslash}p{0.36\linewidth}@{}}
\caption{Representative CPP methods and their applicability to swing-arm troweling.}\label{tab:cpp-comparison}\\
\toprule
Category & Representative method & Map and tool model & Relevance to swing-arm troweling \\
\midrule
\endfirsthead
\toprule
Category & Representative method & Map and tool model & Relevance to swing-arm troweling \\
\midrule
\endhead
Idealized coverage & BCD \cite{choset1998} & Cellular decomposition; fixed footprint & Organizes back-and-forth sweeps within cells, but does not model an active arm or its continuous sweep. \\
Idealized coverage & STC \cite{gabriely2001} & Grid spanning tree; simplified robot & Systematically traverses a grid, but cannot exploit the time-varying tool position created by arm motion. \\
Map-enhanced & Polygon CPP \cite{bahnemann2021,polygoncoverage} & Polygon with holes; fixed tool & Provides alternative sweeps and obstacle-avoiding transfers, but does not represent the arm's instantaneous footprint. \\
Map-enhanced & Fields2Cover \cite{mier2023,fields2cover} & Non-convex field; fixed-width implement & Includes swaths, route ordering, and curvature-constrained connections; the fixed implement model omits arm motion. \\
Map-enhanced & Nobleo FCPP \cite{nobleo} & Occupancy grid; fixed tool disc & Integrates grid coverage with a navigation stack, but does not coordinate arm angle, arm speed, and base speed. \\
Platform-constrained & Mobile-manipulator surface coverage \cite{thakar2022} & Three-dimensional surface; mobile manipulator & Jointly plans base placement and end-tool motion, but does not completely traverse a complex two-dimensional floor. \\
Platform-constrained & SatAOI \cite{lin2025} & Construction obstacle map; base, arm, and disc & Determines workable regions and connectivity, but does not generate a visitation order or continuous coverage trajectory. \\
Proposed method & Hierarchical swing-arm CPP & Two-dimensional obstacle map; time-varying disc sweep & Links rectangular decomposition, visitation, local swinging, and trajectory evaluation; currently assumes a static map and uses heuristic search. \\
\bottomrule
\end{longtable}

\subsection{Summary and Research Gap}

Prior work provides several capabilities needed here, but has not solved complex-map coverage within one unified task. BCD, STC, and related frameworks can represent two-dimensional free spaces containing holes, non-convex boundaries, or multiple obstacles and can decompose regions, generate swaths, or traverse grids. With a tool fixed to the base, they can arrange adjacent sweeps according to a prescribed width. Fields2Cover additionally generates turns and connections for vehicles whose direction and minimum turning radius prevent arbitrary lateral motion. Mobile-manipulator research addresses end-tool reachability and collision relationships, while SatAOI uses the geometry among the base, arm, disc, and obstacles to delimit regions the robot can enter and the disc can treat \cite{lin2025,choset1998,gabriely2001,bahnemann2021,polygoncoverage,mier2023,fields2cover,nobleo,thakar2022}. These abilities remain distributed across different task models and applications: complex-map CPP usually assumes a fixed footprint, detailed robot models are commonly used for local or regular workspaces, and SatAOI terminates at workable-area delimitation. For the specific objective of generating a complete executable route for a swing-arm troweling robot in a known, complex, two-dimensional construction map, three coupled gaps remain.

First, existing CPP methods do not model active arm behavior sufficiently to plan the actual coverage of the end disc. In conventional CPP, the tool is fixed relative to the base, so determining the base route also determines tool position and coverage. For a swing-arm robot, the same base position and heading admit multiple arm poses and therefore different disc positions, swept regions, and whole-robot collision states. Approximating this behavior by the maximum arm reach or an enlarged fixed width confuses geometric reachability in one pose with coverage throughout continuous motion and cannot test the combined collision state of the base, arm, and disc during execution. SatAOI models the geometry and collision relationships among these components but uses them primarily to determine workable positions rather than continuous coverage motion \cite{lin2025}. The present method therefore preserves the relationship among base motion, arm rotation, and the circular-disc sweep during both trajectory generation and evaluation.

Second, complex-map decomposition and local swing motion have not been integrated into a computationally practical planning chain. Existing complex-map CPP frameworks handle polygons with holes, non-convex boundaries, or occupancy grids and produce partitions and global routes, but their local coverage actions typically rely on a fixed tool footprint \cite{bahnemann2021,polygoncoverage,mier2023,fields2cover,nobleo}. Mobile-manipulator methods plan local base and end-tool motions in detail but generally do not decompose and continuously traverse an entire complex two-dimensional floor \cite{thakar2022}. SatAOI identifies workable regions and their connectivity but does not proceed to visitation ordering, inter-region transfers, or coordinated intra-region coverage \cite{lin2025}. Directly enumerating decompositions, visitation orders, and arm trajectories together would cause the number of combinations to grow with the obstacles and work units, undermining computational efficiency and field practicality. The proposed method instead organizes complex-space decomposition, region connectivity, global visitation, and coordinated base--arm coverage as a hierarchical heuristic pipeline that converts map-level decisions into continuous work routes while controlling the candidate-search space.

Third, a static reachability envelope does not guarantee complete coverage during dynamic execution, and prior work does not explicitly model the matching between base linear velocity and arm angular velocity. Maximum arm reach indicates only which locations the disc can touch in some poses. During continuous base advance and periodic arm swinging, coverage continuity depends on the number of sweeps per unit distance. If the base is too fast or the arm too slow, the longitudinal spacing between consecutive sweeps increases, leaving untreated gaps even when the target lies inside the geometric reach envelope. Evaluation based only on that envelope or on an untimed geometric path therefore mistakes ``reachable'' for ``covered.'' This paper incorporates the relationship between base and arm velocities into the coverage analysis and calculates coverage ratio and task time from the area actually swept during continuous motion, exposing the joint effect of speed parameters on gaps and time cost.

In summary, the core research gap is the absence of an end-to-end CPP method that, in a known two-dimensional construction environment with dense obstacles, connects active-arm behavior modeling, global--local coverage planning for a complex map, and continuous swept-area evaluation under coupled base and arm velocities. This work aims to complete the planning chain from a construction map and swing-arm robot model to a complete, continuous, executable coverage trajectory, with final coverage quality determined by the disc sweep under the specified velocity relationship.
